% © 2026 Ing. David Jaroš — CC BY-NC-ND 4.0
%
% This work is licensed under a Creative Commons Attribution-NonCommercial-NoDerivatives
% 4.0 International License (CC BY-NC-ND 4.0).
%
% License History: Earlier drafts (up to v0.3) were released under CC BY 4.0.
% From v0.4 onward, all material is released under CC BY-NC-ND 4.0 to protect
% the integrity of the theoretical work during ongoing academic development.
%
% See LICENSE.md for full license text.
%
% File: research_tracks/alpha_spectral/self_dual_torus_derivation.tex
% Task: derive_or_reject_self_dual_torus_condition
% Priority: CRITICAL
% Date: 2026-05-09
%
% Hard rules:
%   - no fitting
%   - no assuming tau=i
%   - do not use alpha, p=137, eta(i), or B_required
%   - verdict must be PROVED / CONDITIONAL / NO-GO

% UBT-AI-PROVENANCE-BEGIN
% schema: ubt-ai-provenance/v1
% tier: C_working
% ai_assistance: disclosed
% human_review: risk-based
% editorial_responsibility: Ing. David Jaroš
% policy: ../../AI_PROVENANCE.md
% notice: Working material; exhaustive human review is not claimed.
% UBT-AI-PROVENANCE-END

\documentclass[12pt,a4paper]{article}
\usepackage[a4paper, margin=2.5cm]{geometry}
\usepackage{amsmath,amssymb,amsthm}
\usepackage{mathtools}
\usepackage{hyperref}
\usepackage{xcolor}
\usepackage{mdframed}
\usepackage{booktabs}

\newtheorem{theorem}{Theorem}[section]
\newtheorem{proposition}[theorem]{Proposition}
\newtheorem{lemma}[theorem]{Lemma}
\theoremstyle{definition}
\newtheorem{definition}[theorem]{Definition}
\theoremstyle{remark}
\newtheorem{remark}[theorem]{Remark}

\newmdenv[backgroundcolor=green!8, linecolor=green!50!black, linewidth=1.2pt]{resultbox}
\newmdenv[backgroundcolor=yellow!12,linecolor=orange!80!black,linewidth=1.2pt]{gapbox}
\newmdenv[backgroundcolor=red!7,   linecolor=red!55!black,   linewidth=1.2pt]{warnbox}

\newcommand{\PROVED}{\textcolor{green!50!black}{\textbf{[PROVED]}}}
\newcommand{\COND}{\textcolor{blue!70!black}{\textbf{[CONDITIONAL]}}}
\newcommand{\NOGO}{\textcolor{red!80!black}{\textbf{[NO-GO]}}}

\newcommand{\dd}{\mathrm{d}}
\newcommand{\Tr}{\mathrm{Tr}}
\newcommand{\Real}{\mathrm{Re}}

\title{\textbf{Self-Dual Torus Condition from the UBT Action}\[4pt]
\large Derive or Reject $R_t=R_\psi$ Without Input Fitting}
\author{Ing.\ David Jaro\v{s}}
\date{2026-05-09}

\begin{document}
\maketitle

\begin{abstract}
We test whether the UBT action $S[\Theta]$ dynamically selects the square
complex-time torus, i.e.\ $R_t=R_\psi$ (equivalently $\tau=i$ after
normalisation and zero tilt). Starting from the canonical kinetic term,
we compactify Euclidean time $t_E\sim t_E+2\pi R_t$ and imaginary time
$\psi\sim\psi+2\pi R_\psi$, isolate the $(t_E,\psi)$ kinetic sector, and derive
the one-loop effective free energy $F(R_t,R_\psi)$ for independent radii.
Result: the shape mode has a stationary point at $R_t=R_\psi$ and the square
point is locally stable under explicit spectral assumptions, but the action
does not fix the overall size modulus. Therefore the full statement
``UBT selects $R_t=R_\psi$ from first principles'' is conditional.
\end{abstract}

\tableofcontents
\newpage

\section{Starting point: UBT action and Euclidean torus}

The canonical biquaternionic action is
\begin{equation}
S[\Theta]
=\int \dd^4x\,\dd\psi\;\Real\,\Tr\!\left[
\bigl(\nabla^\dagger\Theta\bigr)^*\bigl(\nabla^\dagger\Theta\bigr)-V(\Theta)
\right].
\label{eq:action}
\end{equation}
To define a genuine two-cycle torus we perform the Euclidean rotation
$t\to i t_E$ and compactify both directions:
\begin{equation}
t_E\sim t_E+2\pi R_t,
\qquad
\psi\sim\psi+2\pi R_\psi.
\end{equation}
No relation between $R_t$ and $R_\psi$ is assumed.

\subsection{Isolated $(t_E,\psi)$ kinetic sector}

Keeping only terms relevant to radius selection,
\begin{equation}
S_{t\psi}^{(2)}
=\int_{T^2}\dd t_E\,\dd\psi\int \dd^3x\;
\Tr\!\left[
K_t\,\partial_{t_E}\Theta\,\partial_{t_E}\Theta^\dagger
+K_\psi\,\partial_\psi\Theta\,\partial_\psi\Theta^\dagger
+M^2\Theta\Theta^\dagger
\right],
\label{eq:kinetic-sector}
\end{equation}
with positive kinetic coefficients $K_t,K_\psi>0$ and effective mass parameter
$M^2\ge 0$ (from background and potential terms).

\begin{remark}
If the canonical UBT normalization is isotropic in the two compact cycles, one
has $K_t=K_\psi$.  We keep both coefficients explicit and only impose
$K_t=K_\psi$ when testing the self-dual condition.
\end{remark}

\section{Mode expansion with independent radii}

Expand on the rectangular torus:
\begin{equation}
\Theta(x,t_E,\psi)
=\sum_{m,n\in\mathbb Z}\Theta_{m,n}(x)
\exp\!\left(i\frac{m t_E}{R_t}+i\frac{n\psi}{R_\psi}\right).
\label{eq:mode-expansion}
\end{equation}
Then each mode has quadratic eigenvalue
\begin{equation}
\lambda_{m,n}(R_t,R_\psi)
=K_t\frac{m^2}{R_t^2}+K_\psi\frac{n^2}{R_\psi^2}+M^2.
\label{eq:eigenvalue}
\end{equation}

The one-loop free energy (up to $R$-independent constants and standard
regularisation subtraction) is
\begin{equation}
F(R_t,R_\psi)
=\frac{N_\Theta}{2}\sum_{(m,n)\neq(0,0)}
\log\lambda_{m,n}(R_t,R_\psi),
\label{eq:F-main}
\end{equation}
where $N_\Theta$ is the number of dynamical $\Theta$ components retained in the
sector.

\section{Stationary conditions in shape and scale variables}

Define
\begin{equation}
R_t=s\sqrt{\rho},\qquad R_\psi=\frac{s}{\sqrt{\rho}},
\qquad s>0,\;\rho>0,
\end{equation}
so $\rho=R_t/R_\psi$ (shape) and $s=\sqrt{R_tR_\psi}$ (overall scale).
For isotropic kinetic normalization $K_t=K_\psi=:K$,
\begin{equation}
\lambda_{m,n}(s,\rho)=\frac{K}{s^2}\left(m^2\rho^{-1}+n^2\rho\right)+M^2.
\label{eq:lambda-s-rho}
\end{equation}

\subsection{Shape equation}

\begin{proposition}[Stationarity at square torus]
\label{prop:shape-stationary}
Assume $K_t=K_\psi$ and the mode set is symmetric under $(m,n)\leftrightarrow(n,m)$.
Then
\begin{equation}
\left.\frac{\partial F}{\partial \rho}\right|_{\rho=1}=0,
\end{equation}
so $R_t=R_\psi$ is a stationary point of the shape free energy.
\end{proposition}

\begin{proof}
Under $\rho\mapsto\rho^{-1}$ and index swap $(m,n)\leftrightarrow(n,m)$,
\eqref{eq:lambda-s-rho} is unchanged, hence $F(\rho)=F(\rho^{-1})$.
In variable $y=\ln\rho$, this is $F(y)=F(-y)$, so $F'(0)=0$, i.e.
$\partial_\rho F|_{\rho=1}=0$.
\end{proof}

\subsection{Local stability of the square point}

\begin{proposition}[Positive second variation in the shape mode]
\label{prop:shape-stable}
With assumptions of Proposition~\ref{prop:shape-stationary},
\begin{equation}
\left.\frac{\partial^2 F}{\partial y^2}\right|_{y=0}
=\frac{N_\Theta}{2}\sum_{(m,n)\neq(0,0)}
\frac{4K^2m^2n^2+KM^2(m^2+n^2)}{\left(K(m^2+n^2)+M^2\right)^2}
\ge 0,
\label{eq:second-variation}
\end{equation}
where $y=\ln\rho$. For any nontrivial contributing mode pair the inequality is
strict, so the square torus is a local minimum in the shape direction.
\end{proposition}

\begin{proof}
Differentiate $\log\lambda_{m,n}(y)$ twice at $y=0$ using
$\lambda_{m,n}(y)=Ks^{-2}(m^2e^{-y}+n^2e^{y})+M^2$.
Direct algebra gives the numerator in \eqref{eq:second-variation}; every term is
nonnegative for $K>0$, $M^2\ge0$.
\end{proof}

\subsection{Scale equation}

The scale derivative is
\begin{equation}
\frac{\partial F}{\partial s}
=-\frac{N_\Theta K}{s^3}
\sum_{(m,n)\neq(0,0)}
\frac{m^2\rho^{-1}+n^2\rho}{\lambda_{m,n}(s,\rho)}.
\label{eq:scale-derivative}
\end{equation}
For the pure kinetic+mass sector in \eqref{eq:F-main}, this does not generally
produce an isolated finite critical scale $s_*$ without extra physics
(counterterms, additional potentials, fixed-area constraint, or boundary data).

\begin{gapbox}
\textbf{Key obstruction.}
The current UBT action sector used here constrains the \emph{shape} mode but does
not by itself fix the \emph{scale} modulus $s=\sqrt{R_tR_\psi}$. Therefore
``$R_t=R_\psi$ is selected dynamically'' is not an unconditional consequence of
$S[\Theta]$ alone.
\end{gapbox}

\section{Interpretation of $\tau=i$}

For rectangular torus with zero tilt, the complex modulus is
\begin{equation}
\tau=i\frac{R_\psi}{R_t}=i\rho^{-1}.
\end{equation}
Thus $R_t=R_\psi$ is equivalent to $\tau=i$ in this sector.
From Propositions~\ref{prop:shape-stationary} and~\ref{prop:shape-stable}:

\begin{itemize}
\item Shape stationarity at $\tau=i$ is derived (not assumed).
\item Local stability at $\tau=i$ is derived under explicit spectral assumptions.
\item Full dynamical selection needs one additional ingredient fixing the area
$A_{T^2}=(2\pi)^2R_tR_\psi$ (or equivalent scale condition).
\end{itemize}

\section{Final verdict class}

\begin{resultbox}
\textbf{Verdict: \COND}

From the current UBT kinetic action and spectral one-loop free energy:
\begin{enumerate}
\item The square torus condition $R_t=R_\psi$ (equivalently $\tau=i$ for zero
tilt) is a derived \emph{shape stationary point}.\ \PROVED
\item The square point is locally stable in the shape mode under the explicit
assumptions stated above.\ \PROVED
\item The full pair $(R_t,R_\psi)$ is \emph{not uniquely selected} by this sector,
because no independent equation fixes the overall scale modulus.\ \NOGO
\end{enumerate}
Hence the complete claim ``UBT field equations/selective free energy force
$R_t=R_\psi$'' is currently \textbf{conditional} on an additional scale-fixing
principle.
\end{resultbox}

\begin{warnbox}
\textbf{Hard-rule compliance.}
No fitting used. No assumption $\tau=i$ used. No use of $\alpha$, $p=137$,
$\eta(i)$, or $B_{\mathrm{required}}$. Verdict class is in
\{PROVED, CONDITIONAL, NO-GO\}.
\end{warnbox}

\end{document}
